\section{Literature Review}
Computational notebooks have gained significant relevance in education and research, and recent work has explored how learners interact within these environments. A growing body of research shows that while notebooks support exploratory coding and flexible workflows, they also introduce challenges for understanding the learning process at a fine-grained level. Systems designed to capture cell-level interaction data—such as execution patterns, time on task, and navigation behavior—demonstrate that detailed traces can reveal iterative retries, hesitation, and fragmented workflows that remain invisible in final notebook submissions~\cite{cai2025jupyter}. However, existing approaches of this type are primarily designed for instructor-facing analytics and do not address how such data could support student reflection or metacognitive skill development.

In programming education more broadly, research consistently shows that novice programmers struggle with error identification and debugging, often relying on unsystematic trial-and-error strategies and showing difficulty in locating the source of errors~\cite{robinson2022error}. Additional studies highlight that notebook environments exhibit recurring categories of bugs, suggesting that the structure of notebooks can reinforce certain misunderstandings or error patterns~\cite{de2022bug}. These findings underline the importance of tools capable of making the coding process more transparent and helping learners understand not only the errors they make, but also how they arrive at those errors.

Learning analytics dashboards have been widely explored as mechanisms for supporting student awareness, self-regulation, and instructional decision-making. Systematic reviews show that dashboards can enhance metacognitive insight and provide users with actionable feedback on their learning behavior~\cite{schwendimann2016perceiving}. Despite their potential, most dashboards rely on high-level LMS metrics—such as resource access, assignment completion, or overall activity patterns—which limits their capacity to capture the step-by-step nature of programming work~\cite{paulsen2024learning}. As a result, these systems offer only coarse-grained representations of learning progress.

Research targeting programming courses shows similar limitations. Studies that use notebook-based assignments for analytics tend to track performance at the assignment or submission level, without capturing the underlying cell-level workflow~\cite{sondagar2021elearning}. Other dashboard systems designed to support programming instructors focus on global participation and performance distributions but do not visualize detailed notebook interactions or support student-directed reflection~\cite{groher2024learning}. Together, these approaches illustrate the current focus on high-level, instructor-centered metrics and the scarcity of fine-grained, student-centered analytics for notebook environments.

Across all these lines of research, a consistent gap emerges: existing dashboards rarely provide detailed visualizations of how students work within computational notebooks, nor do they investigate whether such visualizations can improve student reflection, error awareness, or learning outcomes. The predominant emphasis remains on monitoring final outputs, aggregated performance, or instructor dashboards, with limited attention to visualization tools that help learners themselves understand their coding behaviors.

This gap motivates the present project. By developing a student-facing Jupyter dashboard extension that visualizes cell-level interaction histories, execution errors, retry patterns, and timelines of activity, the proposed study aims to contribute empirical evidence on how fine-grained notebook analytics can support metacognitive reflection and reduce repeated programming mistakes. Unlike existing systems, the focus is not solely on teacher monitoring, but on enabling learners to interpret their own behavior, gain awareness of unproductive patterns, and better understand their problem-solving process. The project also complements prior research by evaluating the dashboard through a controlled study, providing insights into both behavioral effects and students' subjective perceptions of usefulness and usability.